\section{Validation Boundaries}
\label{sec:validation-boundaries}

An artificial society is not validated by the realism of a transcript, by one
aggregate match, or by successful execution.  GAWorld separates five claim
levels: individual consistency, macro correspondence, emergent dynamics,
counterfactual response, and reproducibility.  Evidence at one level does not
compensate for absence at another.  This layered view follows established
agent-based-modeling practice in distinguishing implementation verification,
empirical validation, and the intended domain of use
\citep{grimm2006odd,windrum2007empirical,fagiolo2007critical}.

\begin{table}[t]
\centering
\small
\caption{Claim--evidence matrix for the audited GAWorld snapshot.  ``Current
support'' states the strongest defensible claim, not an aspirational goal.}
\label{tab:claim-evidence}
\begin{tabular}{@{}p{0.17\linewidth}p{0.28\linewidth}p{0.25\linewidth}p{0.22\linewidth}@{}}
\toprule
Claim level & Evidence required & Current support & Not established \\
\midrule
Individual consistency & Repeated independent elicitations; longitudinal and perturbation tests; external or held-out criteria & Persistent profiles, memories, diaries, interviews, and one completed two-phase arm are inspectable & Human fidelity, stable identity across models/seeds, or a memory-treatment effect \\
Macro correspondence & Pre-specified external targets; calibrated population; uncertainty and out-of-sample tests & State/economy pipelines can produce aggregates; benchmark fixtures test analysis code & Population representativeness, calibrated macro fit, forecasting, or real-city prediction \\
Emergent dynamics & Multiple seeds and scales; mechanism ablations; network/temporal diagnostics; robust qualitative patterns & Social and network mechanisms emit inspectable state; failures are archived & Polarization, contagion, homophily, or other stable emergent laws \\
Counterfactual response & Replicated randomized arms; verified intervention path; placebo and negative controls; uncertainty & Four paired event/baseline trees demonstrate comparison mechanics & Causal policy effects, sign stability, transportability, or welfare recommendations \\
Reproducibility & Immutable code/config/data/model snapshots; deterministic components; repeated reruns and failure accounting & Hashed result files, seeds/configs in some runs, file-backed traces, and a diagnostic harness & Bitwise replay of hosted LLM calls or complete end-to-end reproduction of every archive \\
\bottomrule
\end{tabular}
\end{table}

\subsection{Individual Consistency}

GAWorld implements the prerequisites for a consistency study: named persistent
agents, structured state, episodic files, recall, reflection, and longitudinal
logs.  These mechanisms allow within-agent questions such as whether a stated
preference survives intervening events or whether an interview answer is
traceable to stored experience.  The current memory case does not answer those
questions comparatively because three of four arms lack phase 2 and the only
complete arm contains one agent.  Generated behavioral coherence is also not
equivalent to correspondence with a real person.  A stronger evaluation would
repeat standardized probes before and after controlled memory perturbations,
separate temporal persistence from prompt sensitivity, and compare against
held-out observations where consent and provenance permit.

\subsection{Macro Correspondence}

Aggregating simulated variables is an implemented capability; calling the
aggregate realistic requires a separate calibration argument.  The audited
economy trace has five agents and one seed, while the benchmark scorecard is a
synthetic diagnostic fixture.  Neither represents a target population or
supports out-of-sample prediction.  Future macro validation should document
the sampling frame, map each simulated construct to an external measurement,
reserve targets for testing rather than calibration, and report uncertainty
across seeds and plausible parameterizations.  Agreement on a few marginal
means would remain insufficient if joint distributions, temporal patterns, or
subgroup errors failed.

\subsection{Emergence and Mechanism}

An emergent claim requires a pattern that is not merely written into a prompt
or initial condition and that recurs under controlled variation.  The present
archive does not meet that standard: the network study records no interactions,
the emotion study has no state files, and the polarization arms have unequal
record counts.  The existence of social-network, exposure, and affect fields
therefore establishes instrumentation, not a social law.  Appropriate tests
would vary population size and topology, ablate candidate mechanisms, quantify
between-seed variability, and trace a macro pattern back to micro-level events.
Where a field remains identically zero, the first question is whether the
mechanism and detector were active, not why society produced no effect.

\subsection{Counterfactual Validity}

The matched-run interface is a useful scaffold for counterfactual experiments,
but the archived policy differences are single contrasts produced by an
LLM-mediated dynamical system.  A common seed does not guarantee common latent
randomness after prompts diverge, and a policy description does not guarantee
that an economic or mobility mechanism changed.  Before interpreting a
difference, researchers should verify the intervention's path from event to
structural state, replicate both arms over seeds, pre-register target metrics,
and include placebo events or negative-control outcomes
\citep{lipsitch2010negative}.  External policy conclusions additionally require
a defensible mapping to the population and institutions of interest.  GAWorld
currently supports constructing and auditing such a study, not replacing it.

\subsection{Reproducibility as a Non-Compensatory Gate}

The artifact hashes used in this paper bind each numerical statement to an
audited byte sequence.  They do not bind the archive to a clean released
runtime: the audit began from a dirty working tree, and hosted providers may
change behavior independently of code, seed, or sampling parameters.  Missing
state files and interrupted arms further show that successful completion must
be recorded per treatment.  Consequently, a result that cannot be reconstructed
from code, configuration, inputs, model identity, and complete outputs should
not be promoted because it looks plausible or matches an external statistic.

For future releases, the minimum reproducibility unit should contain the exact
code revision plus a patch or clean tag, environment lockfile, initial
profiles and maps, provider/model identifiers, prompt and sampling settings,
random seeds, per-arm status, logs, and raw and derived artifacts.  Hosted-model
studies should estimate replay variability rather than promise bitwise
determinism.  The diagnostic benchmark can then serve as a regression gate:
fixtures test parsers and scoring logic, while separately labeled simulation
runs test whether scientific evidence has actually been produced.

\subsection{Intended Use}

At its present evidence level, GAWorld is best used for mechanism prototyping,
trace inspection, hypothesis generation, controlled software experiments, and
the design of more rigorous artificial-society studies.  The archive does not
support deploying GAWorld as a substitute population, forecasting real social
outcomes, ranking policies, or making decisions about identifiable people.
Those are not merely stronger wordings of the same result; they are different
claims requiring external data, identification, uncertainty quantification,
and governance that the current artifacts do not provide.
